%% REPLACE SXXXXXX with your student number
\def\studentNumber{S2149988}


%% START of YOUR ANSWERS
%% Add answers to the questions below, by replacing the text inside the brackets {} for \youranswer{ "Text to be replaced with your answer." }. 
%
% Do not delete the commands for adding figures and tables. Instead fill in the missing values with your experiment results, and replace the images with your own respective figures.
%
% You can generally delete the placeholder text, such as for example the text "Question Figure 2 - Replace the images ..." 
%
% There are 18 TEXT QUESTIONS (a few of the short first ones have their answers added to both the Introduction and the Abstract). Replace the text inside the brackets of the command \youranswer with your answer to the question.
%
% There are also 3 "questions" to replace some placeholder FIGURES with your own, and 3 "questions" asking you to fill in the missing entries in the TABLES provided. 
%
% NOTE! that questions are ordered by the order of appearance of their answers in the text, and not by the order you should tackle them. Specifically, you cannot answer Questions 2, 3, and 4 before concluding all of the relevant experiments and analysis. Similarly, you should fill in the TABLES and FIGURES before discussing the results presented there. 
%
% NOTE! If for some reason you do not manage to produce results for some FIGURES and TABLES, then you can get partial marks by discussing your expectations of the results in the relevant TEXT QUESTIONS (for example Question 8 makes use of Table 1 and Figure 2).
%
% Please refer to the coursework specification for more details.


%% - - - - - - - - - - - - TEXT QUESTIONS - - - - - - - - - - - - 



%% Question 1:
\newcommand{\questionOne} {
\youranswer{training accuracy increases over all the epochs displayed, whereas, the validation accuracy reaches a peak
around epoch 12 (~83\%) before steadily declining. We see that in figure 1b, the training error consistently decreases
across all the epochs shown, whereas the validation error reaches a minimum around epoch 9 before increasingly at a
steady rate. Additionally, we can see that the generalisation gap increases significiantly from epoch 9 onwards.
This, widening gap between the errors shows that the model is significiantly overfitting the data.
}
}

%% Question 2:
\newcommand{\questionTwo} {
\youranswer{
}
%    Question 2 - Present your network width experiment results by using the relevant figure and table}
}


%%

%% "Question 2:
%Present your network width experiment results by using the relevant figure and table.
%[see also Question 3]"
%Correctly states which models have been overfit...	"... and points to evidence in the figure for the specific training instance/models that overfit.
%
%[This can be a statement on which models have their generalisation gap increasing or pointing out specific values on the figures].
%
%[check Table 1 and Figure 2  that their statements are mostly concistent with their results]."
%
%"Question 3
%[Marking 2 & 3 together, doesn't matter which answer has the info]:
%Discuss whether varying width affects the results in a consistent way, and whether the results are expected and match well with the prior knowledge (by which we mean your expectations as are formed from the relevant Theory and literature)"
%"Correctly discusses the trend (or lack thereof) of the generalisation gap/overfitting as hyperparamer values change.
%
%[check Table 1 and Figure 2 in case something else holds in the students case; what matters here is consistency with their own results]."	Points out the model with the best validation accuracy (from this series).	"Points out the model that overfits less (from this series).
%
%[these are not necessarily the same, so the student should not just say ""this model is best"". If they simply state this charitably assume they are talking about overfitting]."	Uses relevant literature to form their expectations of how the experiment will/should have or would have been expected to turn out.	Forms expectations via (reasonably) arguing based on model capacity.	Goes beyond these expectations in successfully addressing Questions 2&3 (not simply added information but either things that add to the argument or an excellent delivery of previous points).

%% Question 3:
\newcommand{\questionThree} {
\youranswer{models __ and __ have been overfit, this can be seen in figure 2b as the generalisation gap for these models
increases ___ over the epochs, whereas model ___ ___ do not display such generalisation gaps.
This suggests that as the width of a model increases the greater the risk of overfitting. This

increases significiantly whereas models __ and __ ____.
is the case as it's generalisation gap increases


    for the widths tested, increasing the network's width consistently improves validation accuracy and reduces
both training and validation errors. This suggests that wider networks improve overall performance.
This behaviour is typical up to a point, as wider networks have greater capacity to model complex patterns.
However, it is also known that increasing a network's width beyond a certain point, results in either overfitting or reducing performance gains.
We can see the early signs of overfitting within the width 128 as it's generalisation gap larger than for 64 and 32.
%    Question 3 - Discuss whether varying width affects the results in a consistent way, and whether the results are expected and match well with the prior knowledge (by which we mean your expectations as are formed from the relevant Theory and literature)}
}
}

%% Question 4:
\newcommand{\questionFour} {
\youranswer{ }
%    Question 4 - Present your network depth experiment results by using the relevant figure and table}
}

%% Question 5:
\newcommand{\questionFive} {
\youranswer{
%    Question 5 - Discuss whether varying depth affects the results in a consistent way, and whether the results are expected and match well with the prior knowledge (by which we mean your expectations as are formed from the relevant Theory and literature)}
for the depths tested, increasing the network depth from one to two hidden layers yields a small improvement in validation accuracy, but extending the depth further to three layers provides no additional benefit. Although the training error continues to decrease with depth, the validation error plateaus, indicating that additional capacity is not translating into better generalisation.
The slight rise in the generalisation gap from depth two to depth three suggests the earliest onset of overfitting: the model is becoming better at fitting the training data without achieving corresponding gains on unseen data.
This behaviour is consistent with standard theory, as depth two multi layer perceptrons are known to be universal function approximators, capable of representing any continuous function given sufficient width.
}
}


%% Question 6:
%    Question 6 - Explain how one could use a combination of L1 and L2 regularisation. Discuss any potential benefits of this approach. Present an equation for this combined loss function and explain all relevant variables and hyperparameters.}
\newcommand{\questionSix} {
\youranswer{
We can combine L1 and L2 regularisation into a single regulariser by introducing a separate parameter for each
regularising term. L1 encourages sparsity by penalising the absolute magnitude of the weights,
while L2 smooths the weights by penalising large weights more strongly through their squared magnitude.

This combined regulariser could provide a flexible balance between weight sparsity and weight smoothness,
whilst reducing the negative consequences of each individual regulariser. For example:
L2 can prevent L1 from discarding weaker correlated weights making the model more stable,
whilst L1 provides the ability to eliminate truly irrelevant weights that L2 may maintain.

We can define this loss function with combined regulariser using the equation below:

\[
L_{\text{total}}
= L_{\text{data}}
+ \lambda_{1}\|w\|_{1}
+ \lambda_{2}\|w\|_{2}^{2},
\]

\begin{itemize}
    \item \( L_{\text{total}} \): The total loss the optimiser is minimising.
    \item \( L_{\text{data}} \): The loss produced from the data (e.g.\ cross-entropy or MSE).
    \item \( w \): Vector of the model weights.
    \item \( \|w\|_{1} \): L1 term, promoting sparsity.
    \item \( \|w\|_{2}^{2} \): L2 term, encourages smoother weights.
    \item \( \lambda_{1} \): Strength of the L1 regularisation term.
    \item \( \lambda_{2} \): Strength of the L2 regularisation term.
\end{itemize}



}
}



%% Question 7:
%    Question 7 - Assume you had a limited experiment budget for exactly \textbf{4} experiment runs dedicated to testing the effectiveness of your proposed combination of L1 and L2 Regularization in Section~\ref{sec:L1L2combo}. Argue for which 4 configurations to run based on any previous results or other information you consider relevant. Run those experiments and add the results in Table~4 and Figure~5 (you will then discuss those results alongside all others in the next question below).
\newcommand{\questionSeven} {
\youranswer{Due to limited experimental budget of four runs, we focused on the interaction between L1 and L2 rather than
just finding a single “best” setting. From our earlier tests, both L1 and L2 have useful ranges around 1e-13–1e-9,
with performance peaking near 1e-11 and degrading when too large.
Guided by that, we probed two balanced points and two asymmetric points to test whether gains are driven primarily by one term or if a complementary effect emerges.

    Specifically, we ran:
    \begin{itemize}
        \item Balanced small: \(\lambda_1=1\text{e-}11,\ \lambda_2=1\text{e-}11\) — tests a minimal joint penalty near each method’s individual sweet spot.
        \item Balanced moderate: \(\lambda_1=1\text{e-}09,\ \lambda_2=1\text{e-}09\) — tests whether increasing both together maintains or improves generalisation.
        \item Stronger L1: \(\lambda_1=1\text{e-}09,\ \lambda_2=1\text{e-}11\) — tests if L1 drives most of the benefit with only a small L2 for smoothing.
        \item Stronger L2: \(\lambda_1=1\text{e-}11,\ \lambda_2=1\text{e-}09\) — tests if additional L2 helps when L1 is kept small.
    \end{itemize}

    This set stays within the four-run budget while answering a clearer question: do we see synergy between L1 and L2, or does one term dominate? Using balanced vs asymmetric pairs lets us attribute any improvement to either interaction effects or to a single regulariser, and the chosen magnitudes reflect the effective region indicated by our prior results.
}
}


%% Question 8:
\newcommand{\questionEight} {
\youranswer{
%    Question 8 - Explain the experimental details (e.g. hyperparameters), discuss the results in terms of their generalisation performance and overfitting. Select and test the best performing model as part of this analysis.
Comparing our results to the baseline in our previous experiments, shows that for the parameters we have chosen,
    all dual combinations of L1L2 regularisation result in reduced performance.
    Additionally, if we compare their performance to baseline L1 and L2 it reveals that for these parameters any combination is
    both worse than without regularisation, but also worse than using them individually.
    From the parameters we have tested, we can see that a more
    combining them
    From the parameters we have tested, we can see that a more moderate L1 provides better accuracy, compared to L2 which provides minimal gains with a moderate L1 and in fact hinders performance if L1 is minimal.
    We can see that L1, provides a larger improvement for performance.

If we use our

The results show that a larger L1 provides the majority
whilst L2 provides minimal improvements within regularisation.
In fact, using a moderate L2 without a moderate L1 reduces performance compared to the combined small version. A
the improved
}
}


%% Question 9:
\newcommand{\questionNine} {
\youranswer{
%    Question 9 - A colleague proposed using a custom activation function (which can be identified as CustomActivationLayer() in your Coursework\_1.ipynb notebook and the mlp package). Run the prepared experiment in the "Problematic Training" cell in your Coursework\_1.ipynb notebook, and display the results in Figure~6. Analyse the results. Is the model overfitting? Explain the cause of the problematic performance of the model?

The experiment using the proposed custom activation function $( f(x) = \frac{1}{10(1 + e^{-x})} )$ resulted in severe underfitting.
Both training and validation accuracies remained around 0.02, which is effectively random performance for a 47-class classification task (( 1/47 ~ 0.0213 )).
This indicates that the model failed to learn any meaningful representations during training.
The poor performance is explained by the nature of the activation function.
It is a scaled sigmoid, which limits outputs to the narrow range ((0, 0.1)).
This scaling compresses activations so strongly that the gradients become extremely small,
causing vanishing gradients early in backpropagation.
As a result, the network’s weights receive negligible updates, preventing effective learning.
This effect can be clearly seen in the gradient plot, where gradients quickly decay toward zero before reaching deeper layers.
In contrast, we repeated the experiment swapping back to a standard ReLU activation
which maintained non-zero gradients throughout the network, enabling learning.
}
}




%% Question 10:
\newcommand{\questionTen} {
\youranswer{Other regularisation methods such as L1, L2 and Dropout all provided some improvements, but none were as effective as label smoothing.
    Combining L1 and L2 regularisation did not yield better results than using them individually,
    suggesting that their benefits do not compound effectively for our specific task and architecture.
    Finally, we found that a scaled sigmoid function is not a suitable regularisation technique, as it caused vanishing gradients and prevented effective learning.
%    Question 10 - Briefly draw your conclusions based on the results from the previous sections (what are the take-away messages?), discussing them in the context of the overall literature, and conclude your report with a recommendation for future directions}
}
}

%% - - - - - - - - - - - - FIGURES - - - - - - - - - - - - 

%% Question Figure 2:
\newcommand{\questionFigureTwo} {
\youranswer{%    Question Figure 2 - Replace the images in Figure 2 with figures depicting the accuracy and error, training and validation curves for your experiments varying the number of hidden units.
\begin{figure}[t]
    \centering
    \begin{subfigure}{\linewidth}
        \includegraphics[width=\linewidth]{figures/width_acc}
        \caption{accuracy by epoch}
        \label{fig:width_acccurves}
    \end{subfigure} 
    \begin{subfigure}{\linewidth}
        \centering
        \includegraphics[width=\linewidth]{figures/width_error}
        \caption{error by epoch}
        \label{fig:width_errorcurves}
    \end{subfigure} 
    \caption{Training and validation curves in terms of classification accuracy (a) and cross-entropy error (b) on the EMNIST dataset for different network widths.}
    \label{fig:width}
\end{figure} 
}
}

%% Question Figure 3:
\newcommand{\questionFigureThree} {
\youranswer{}
%    Question Figure 3 - Replace these images with figures depicting the accuracy and error, training and validation curves for your experiments varying the number of hidden layers.
%
\begin{figure}[t]
    \centering
    \begin{subfigure}{\linewidth}
        \includegraphics[width=\linewidth]{figures/depth_acc}
        \caption{accuracy by epoch}
        \label{fig:depth_acccurves}
    \end{subfigure} 
    \begin{subfigure}{\linewidth}
        \centering
        \includegraphics[width=\linewidth]{figures/depth_error}
        \caption{error by epoch}
        \label{fig:depth_errorcurves}
    \end{subfigure} 
    \caption{Training and validation curves in terms of classification accuracy (a) and cross-entropy error (b) on the EMNIST dataset for different network depths.}
    \label{fig:depth}
\end{figure} 
}

%% Question Figure 4:
\newcommand{\questionFigureFour} {
\youranswer{
%    Question Figure 4 - Replace these images with figures depicting the Validation Accuracy and Generalisation Gap (difference between validation and training error) for each of the experiment results varying the Dropout inclusion rate, and L1/L2 weight penalty depicted in Table 3 (including any results you have filled in).
\begin{figure*}[t]
    \centering
    \begin{subfigure}{.475\linewidth}
        \includegraphics[width=\linewidth]{figures/dropout_plot}
        \caption{Accuracy and error by inclusion probability.}
        \label{fig:dropoutrates}
    \end{subfigure} 
    \begin{subfigure}{.475\linewidth}
        \centering
        \includegraphics[width=\linewidth]{figures/wd_plot}
        \caption{Accuracy and error by weight penalty.}
        \label{fig:weightrates}
    \end{subfigure} 
    \caption{Accuracy and error by regularisation strength of each method (Dropout and L1/L2 Regularisation).}
    \label{fig:hp_search}
\end{figure*}
}
}

%% Question Figure 5:
\newcommand{\questionFigureFive} {
\youranswer{
%    Question Figure 5 - Add figures depicting the Validation Accuracy and Generalisation Gap (difference between validation and training error) for each of your 4 experiments  varying your hyperparameters for the combination of L1 and L2 regularisation (based on your implementation) depicted in Table 4.
\begin{figure*}[t]
    \centering
    \begin{subfigure}{.475\linewidth}
        \includegraphics[width=\linewidth]{figures/l1_1e-11_l2_1e-11_acc_gen_gap}
        \caption{Balance Small L1L2 Regularisation}
        \label{fig:balancedsmallreg}
    \end{subfigure}
    \begin{subfigure}{.475\linewidth}
        \includegraphics[width=\linewidth]{figures/l1_1e-09_l2_1e-09_acc_gen_gap}
        \caption{Balance Small L1L2 Regularisation}
        \label{fig:balancedmoderatereg}
    \end{subfigure}
    \begin{subfigure}{.475\linewidth}
        \includegraphics[width=\linewidth]{figures/l1_1e-09_l2_1e-11_acc_gen_gap}
        \caption{Stronger L1 Regularisation}
        \label{fig:strongerl1}
    \end{subfigure}
    \begin{subfigure}{.475\linewidth}
        \includegraphics[width=\linewidth]{figures/l1_1e-11_l2_1e-09_acc_gen_gap}
        \caption{Stronger L2 Regularisation}
        \label{fig:strongerl2}
    \end{subfigure}
    \label{fig:l1l2regs}
\end{figure*}
}
}

%% Question Figure 6:
\newcommand{\questionFigureSix} {

\youranswer{
%    Question Figure 6 - Add figures depicting the Training Accuracy, Error, and Gradient Flow Across Layers for the 1 experiment run with the Custom Activation function (produce these by running the "Problematic Training" cell in the Coursework\_1.ipynb notebook).
\begin{figure}[t]
    \centering
    \begin{subfigure}{\linewidth}
        \includegraphics[width=\linewidth]{figures/acc_custom_activation_curve.png}
        \caption{accuracy by epoch}
        \label{fig:custom_act_acccurves}
    \end{subfigure}
    \begin{subfigure}{\linewidth}
        \centering
        \includegraphics[width=\linewidth]{figures/error_custom_activation_curve}
        \caption{error by epoch}
        \label{fig:custom_act_errorcurves}
    \end{subfigure}
    \begin{subfigure}{\linewidth}
        \centering
        \includegraphics[width=\linewidth]{figures/custom_activation_grad_flow}
        \caption{gradient flow across layers}
        \label{fig:custom_act_gradflow}
    \end{subfigure}
    \caption{Training accuracy, error and gradient flow for a custom activation function}
    \label{fig:custom_act}
\end{figure}

}
}

%% - - - - - - - - - - - - TABLES - - - - - - - - - - - - 

%% Question Table 1:
\newcommand{\questionTableOne} {
\youranswer{
\begin{table}[t]
    \centering
    \begin{tabular}{c|ccc}
    \toprule
        \# Hidden Units & Val. Acc. & Train Error & Val. Error\\
    \midrule
    32 & 0.751 & 0.788 & 0.819 \\
    64 & 0.799 & 0.579 & 0.637 \\
    128 & 0.830 & 0.420 & 0.522 \\
    \bottomrule
    \end{tabular}
    \caption{Validation accuracy (\%) and training/validation error (in terms of cross-entropy error) for varying network widths on the EMNIST dataset.}
    \label{tab:width_exp}
\end{table}
}
}

%% Question Table 2:
\newcommand{\questionTableTwo} {
\youranswer{
%Question Table 2 - Fill in Table 2 with the results from your experiments varying the number of hidden layers.
%
\begin{table}[t]
    \centering
    \begin{tabular}{c|ccc}
    \toprule
        \# Hidden Layers & Val. Acc. & Train Error & Val. Error \\
    \midrule
    1           &   0.829   &  0.426 &   0.524   \\
    2           &   0.832   &  0.406 &   0.505   \\
    3           &   0.831   &  0.403 &   0.505   \\
    \bottomrule
    \end{tabular}
    \caption{Validation accuracy (\%) and training/validation error (in terms of cross-entropy error) for varying network depths on the EMNIST dataset.}
    \label{tab:depth_exps}
\end{table}
}
}

%% Question Table 3:
\newcommand{\questionTableThree} {
\youranswer{
%Question Table 3 - Fill in Table 3 with the results from your experiments for the missing hyperparameter values for each of L1 regularisation, L2 regularisation, Dropout and label smoothing (use the values shown on the table).
%
\begin{table*}[t]
    \centering
    \begin{tabular}{c|c|ccc}
    \toprule
        Model    &  Hyperparameter value(s) & Validation accuracy & Train Error & Validation Error \\
    \midrule
    \midrule
        Baseline &  -         &   0.757   &  0.783 &   0.801   \\
    \midrule
        \multirow{4}*{Dropout}
                & 0.35        &   0.023   &  3.782 &   3.783   \\
                & 0.55        &   0.641   &  2.987 &   2.993   \\
                & 0.75        &   0.716   &  1.625 &   1.634   \\
                & 0.95        &   0.760   &  0.809 &   0.828   \\
    \midrule
        \multirow{4}*{L1 penalty}
                &  1e-13       &   0.755   &  0.778 &   0.798   \\
                &  1e-11       &   0.768   &  0.748 &   0.766   \\
                &  1e-09       &   0.765   &  0.766 &   0.782   \\
                &  1e-07       &   0.758   &  0.790 &   0.806   \\
                &  1e-05       &   0.038   &  3.754 &   3.754   \\
    \midrule
        \multirow{4}*{L2 penalty}  
                &  1e-13       &   0.765   &  0.758 &   0.780   \\
                &  1e-11       &   0.769   &  0.745 &   0.764   \\
                &  1e-09       &   0.762   &  0.774 &   0.799   \\
                &  1e-07       &   0.756   &  0.790 &   0.813   \\
                &  1e-05       &   0.437   &  2.203 &   2.218   \\
    \midrule
            Label smoothing & 0.1 & 0.780   &  1.431 &   1.448   \\
    \bottomrule
    \end{tabular}
    \caption{Results of all hyperparameter search experiments. \emph{italics} indicate the best results per series (Dropout, L1 Regularisation, L2 Regularisation, Label smoothing) and \textbf{bold} indicates the best overall.}
    \label{tab:hp_search}
\end{table*}
}
}


%% Question Table 4:
\newcommand{\questionTableFour} {
\youranswer{
Question Table 4 - Add a table with the results from your 4 experiments (by which we mean: 4 training runs, NOT 4 sets of experiments) varying your hyperparameters for the combination of L1 and L2 regularisation (based on your implementation).

\begin{table*}[t]
    \centering
    \begin{tabular}{cc|ccc}
    \toprule
    \ L1 Coefficient & L2 Coefficient & Val. Acc. & Train Error & Val. Error \\
    \midrule
    1e-11       & 1e-11       &   0.747   &  0.828 &   0.845   \\
    1e-11       & 1e-09       &   0.744   &  0.847 &   0.866   \\
    1e-09       & 1e-11       &   0.750   &  0.823 &   0.836   \\
    1e-09       & 1e-09       &   0.751   &  0.794 &   0.818   \\
    \bottomrule
    \end{tabular}
\end{table*}
}
}

%% END of YOUR ANSWERS
